
\documentclass[11pt,a4paper]{article}

\usepackage[margin=1in]{geometry}
\usepackage{amsmath,amssymb,amsthm}

\title{Results on Diana Thesis}
\author{}
\date{}

\begin{document}
\maketitle

\section{Revisiting Proposition 3.5}

Diana's Proposition 3.5 describes which simple modules
occur in the last nonzero socle layer of an indecomposable
injective $\mathfrak{sp}_{\infty}$-module $\Gamma_\lambda$,
where $\lambda$ is a partition.

The simple modules are denoted by
$\Gamma_{\langle\mu\rangle}$, with $\mu$ another partition.

From Proposition 3.4, the last nonzero socle layer
occurs at index $k_{\max}+1$, where
\[
k_{\max}=\lambda_2+\lambda_4+\lambda_6+\cdots.
\]

The unique maximal subpartition of $\lambda$
of the form $(2\gamma)^T$ is
\[
\eta_{\max}
=(\lambda_2,\lambda_2,\lambda_4,\lambda_4,\ldots).
\]

Here $T$ denotes transposition of Young diagrams,
and $2\gamma$ is obtained by doubling each part
of the partition $\gamma$.

Let $d$ be the partition obtained by rearranging
the nonzero differences
\[
\lambda_1-\lambda_2,\quad
\lambda_3-\lambda_4,\quad
\lambda_5-\lambda_6,\quad\ldots
\]
in decreasing order.

Diana claims that if
\[
\mu\geq_{\mathrm{lex}}d,\qquad
|\mu|=|d|,\qquad
\ell(\mu)\leq\ell(d),
\]
then $\Gamma_{\langle\mu\rangle}$ occurs
in the final socle layer.

Here $|\mu|$ denotes the number of boxes in the
Young diagram of $\mu$, while $\ell(\mu)$ denotes
the number of its nonzero parts.

\subsection*{First observation: lexicographic versus dominance}

The first issue is the use of lexicographic order
rather than dominance order.

Lexicographic order compares partitions at their
first differing entry. Dominance order instead
requires all partial sums to satisfy
\[
\sum_{i=1}^{j}\mu_i
\geq
\sum_{i=1}^{j}d_i
\qquad\text{for every }j.
\]

Dominance implies lexicographic order, but the
converse is false.

For example, consider
\[
\lambda=(9,5,5,1,1),\qquad
d=(4,4,1),\qquad
\mu=(5,2,2).
\]

Although $\mu$ is lexicographically greater
than $d$, the dominance condition fails:
\[
\mu_1+\mu_2=7<8=d_1+d_2.
\]

In fact, the corresponding Littlewood--Richardson
coefficient is
\[
N^{(9,5,5,1,1)}_{(5,5,1,1),(5,2,2)}=0.
\]

Therefore, the proposition is not correct under
the lexicographic order explicitly defined
in Diana's thesis.

It remains possible that dominance order was
intended, since the cited reference discusses
dominance. However, the proof itself also
uses lexicographic order.

\subsection*{Second observation: the tableau construction}

A separate issue concerns the construction used
in the proof.

The proposed algorithm begins in the lowest
incomplete row of the skew diagram, placing
the largest unused number $x$, followed by
$x-1,x-2,\ldots,1$ in successive incomplete
rows above.

The process is repeated until all boxes are
filled, after which the entries in each row
are rearranged into increasing order.

However, this procedure does not always work,
even when the dominance condition is satisfied.

For example, take
\[
\lambda=(4,2,2,1,1),\qquad
d=(2,1,1),\qquad
\mu=(2,2).
\]

The partition $\mu$ dominates $d$, and
\[
N^{(4,2,2,1,1)}_{(2,2,1,1),(2,2)}=1.
\]

A valid LR tableau exists with two 1's in the
highest skew row and one 2 in each of the
two lower skew rows.

Diana's algorithm instead starts by placing
a 2 in the lowest row and a 1 in the middle row.

This prevents a valid completion because
both 1's must occupy the highest row to
satisfy the LR lattice-word condition.

Thus, replacing lexicographic order with
dominance order gives the correct statement,
but the proposed tableau construction still
requires a different argument.

\section{Refinement Using Kostka Numbers}

The previous observations suggest a more direct
way to characterize the entire final socle layer.

Instead of determining only which simple modules
occur, the multiplicity of each constituent
can be described explicitly using Kostka numbers.

By Proposition 3.4, the multiplicity of
$\Gamma_{\langle\mu\rangle}$ in the final
socle layer is
\[
N^\lambda_{\eta_{\max},\mu}.
\]

The interesting observation is that the skew
diagram $\lambda/\eta_{\max}$ consists of
horizontal rows occupying disjoint columns.

The lengths of these rows are precisely the
parts of $d$.

Consequently, the corresponding skew Schur
function factors as
\[
s_{\lambda/\eta_{\max}}
=
h_{d_1}h_{d_2}\cdots h_{d_{\ell(d)}}.
\]

Here $s_{\lambda/\eta_{\max}}$ denotes a
skew Schur function, and $h_r$ denotes the
complete homogeneous symmetric function
of degree $r$.

Young's rule gives
\[
h_{d_1}h_{d_2}\cdots h_{d_{\ell(d)}}
=
\sum_{\mu\vdash |d|}K_{\mu,d}s_\mu,
\]
where $s_\mu$ is the Schur function indexed
by $\mu$, and the sum ranges over all
partitions $\mu$ of $|d|$.

The coefficients $K_{\mu,d}$ are Kostka numbers.
They count semistandard Young tableaux of
shape $\mu$ and weight $d$.

On the other hand, the skew Schur function
also satisfies
\[
s_{\lambda/\eta_{\max}}
=
\sum_\mu
N^\lambda_{\eta_{\max},\mu}s_\mu.
\]

Comparing coefficients in these two expansions
gives
\[
\boxed{
N^\lambda_{\eta_{\max},\mu}=K_{\mu,d}.
}
\]

Therefore, the entire last nonzero socle layer
can be expressed as
\[
\boxed{
\operatorname{soc}_{k_{\max}+1}\Gamma_\lambda
=
\bigoplus_{\mu\vdash |d|}
K_{\mu,d}\Gamma_{\langle\mu\rangle}.
}
\]

Since Kostka numbers are positive precisely
when $\mu$ dominates $d$, the result also
provides an exact characterization of
which simple modules occur.

This gives a proof of the dominance-order
version of Proposition 3.5 without relying
on the original tableau construction.

\section{Some Immediate Consequences}

\subsection*{Simple final socle layers}

Diana already characterizes when the
final socle layer is simple.

The Kostka description provides another
way to see this.

If $d$ has only one nonzero part, its
complete homogeneous symmetric function
is itself a Schur function.

Conversely, if $d$ has at least two
nonzero parts, more than one distinct
Schur function occurs in the expansion.

Hence
\[
\boxed{
\text{The final socle layer is simple}
\iff \ell(d)\leq1.
}
\]

This recovers Diana's Corollary 3.8,
including the empty-partition case.

\subsection*{Multiplicity-free final socle layers}

A natural extension is to characterize
when all simple constituents occur
with multiplicity one.

Such a decomposition is called
multiplicity-free.

Suppose first that $d$ has two
nonzero parts, so $d=(a,b)$.

Pieri's rule gives
\[
h_ah_b
=
\sum_{j=0}^{\min(a,b)}
s_{(a+b-j,j)}.
\]

Every Schur function appears with
coefficient one, so the final
socle layer is multiplicity-free.

Now suppose $d$ has $m\geq3$
nonzero parts and let $n=|d|$.

A standard formula for Kostka numbers gives
\[
K_{(n-1,1),d}=m-1.
\]

Since $m\geq3$, this multiplicity is
at least two.

Therefore, some simple constituent
necessarily occurs more than once.

This gives the characterization
\[
\boxed{
\text{The final socle layer is multiplicity-free}
\iff \ell(d)\leq2.
}
\]

This extends Diana's classification of
simple final socle layers to all
multiplicity-free final socle layers.

\end{document}
